\documentclass[12pt]{report}
\usepackage[utf8]{inputenc}
\usepackage[T1]{fontenc}
\usepackage[spanish,shorthands=off]{babel}
\usepackage{geometry}
\usepackage{setspace}
\usepackage{xcolor}
\usepackage{hyperref}
\hypersetup{
    colorlinks=true,
    linkcolor=blue,
    citecolor=blue,
    filecolor=magenta,
    urlcolor=blue,
}
\usepackage{graphicx}
\usepackage{tikz}
\usetikzlibrary{arrows.meta,positioning,shapes.geometric,fit,calc,babel}
\newenvironment{tikzspanish}[1][]{%
  \begingroup
  \catcode`\;=12\relax
  \catcode`\:=12\relax
  \catcode`\!=12\relax
  \catcode`\?=12\relax
  \catcode`\< =12\relax
  \catcode`\>=12\relax
  \begin{tikzpicture}[#1]}{%
  \end{tikzpicture}\endgroup}
\usepackage{titlesec}
\usepackage{float}
\usepackage{verbatim}

\geometry{a4paper, margin=1in}
\onehalfspacing

\title{\textbf{Whitepaper Red Biet}\\
Distribución Justa de Recursos Sociales Mediante Blockchain}
\author{Team Red Biet -- MinTIC Bootcamp 2025}
\date{2025}

\begin{document}
\maketitle

\begin{abstract}
Red Biet (BietNetwork) es un marco bioeconómico descentralizado orientado a transformar la distribución de recursos sociales en comunidades vulnerables. 
El modelo se basa en unidades vivas llamadas \textit{Biets}, gobernadas por una DAO, con identidad digital descentralizada (DID), trazabilidad en blockchain y un token nativo de utilidad y gobernanza denominado Biet Coin (BGT). 
Este documento unifica la visión conceptual, el estado del arte, la arquitectura Web3, los modelos criptográficos, los flujos de dinero y operación, así como los indicadores de impacto social y mecanismos de transparencia, en el contexto de la opción de proyecto ``Distribución justa de los recursos sociales mediante Blockchain''.
\end{abstract}

\tableofcontents
\clearpage

%----------------------------------
\chapter{Resumen ejecutivo}
%----------------------------------

Los modelos de asistencia social tradicionales en América Latina se enfrentan a problemas estructurales de ineficiencia, opacidad y falta de sostenibilidad. 
En muchos casos, los subsidios y transferencias condicionadas no logran romper los ciclos de pobreza, ni fortalecen capacidades productivas ni la autonomía de las comunidades [1], [2].

Red Biet propone una alternativa: un sistema en el que cada unidad social (Biet) se convierte en una célula viva productiva, conectada a una red blockchain que garantiza trazabilidad, participación y distribución justa de recursos. 
A través de identidad digital descentralizada (DID), gobernanza mediante DAO, contratos inteligentes en una red EVM y el token Biet Coin (BGT), el proyecto articula una infraestructura técnico-social para canalizar recursos de manera transparente, medir impacto y fortalecer economías locales.

El presente whitepaper desarrolla:
\begin{itemize}
    \item El contexto y el problema de la asistencia tradicional.
    \item Un estado del arte de blockchain para inclusión social, DID y DAOs.
    \item El diseño conceptual de Red Biet y el rol de BGT.
    \item La arquitectura Web3 e integración con Next.js y EVM.
    \item Flujos operativos, financieros, de gobernanza y auditoría.
    \item Un anexo técnico con esbozos de contratos inteligentes.
\end{itemize}

%----------------------------------
\chapter{Introducción y estado del arte}
%----------------------------------

\section{Contexto del problema}
Los programas sociales en la región suelen operar con estructuras centralizadas, registros fragmentados y sistemas de información cerrados. 
Diversos estudios de la OCDE y el Banco Mundial señalan que el gasto social no siempre se traduce en mejoras sostenibles, en parte por la falta de mecanismos robustos de seguimiento, evaluación y control ciudadano [1], [2].

La corrupción, la captura de rentas y la mala focalización de subsidios afectan la legitimidad de las políticas públicas. 
Informes de Transparencia Internacional muestran cómo la percepción de corrupción en el sector público sigue siendo elevada en buena parte de América Latina [3]. 
Al mismo tiempo, los sistemas de protección social aún dependen de bases de datos cerradas, sin identidad digital interoperable ni trazabilidad abierta.

\section{Estado del arte: blockchain, identidad y gobernanza}
En la última década han surgido propuestas que integran blockchain con:
\begin{itemize}
    \item \textbf{Identidad digital descentralizada (DID):} la recomendación del W3C sobre DID define un estándar para identidades auto-soberanas, controladas por los propios usuarios, verificables criptográficamente y portables entre sistemas [4]. Ver: \href{https://www.w3.org/TR/did-core/}{w3.org/TR/did-core/}.
    \item \textbf{Gobernanza descentralizada (DAOs):} autores como Buterin han descrito las DAOs como estructuras de gobernanza programables donde reglas y decisiones se ejecutan mediante contratos inteligentes [5].
    \item \textbf{Gestión de cadenas de suministro y trazabilidad:} Cong y He analizan cómo blockchain permite mejorar la transparencia en cadenas de valor complejas, reduciendo la necesidad de intermediarios de confianza [6].
\end{itemize}

En paralelo, la Comisión Europea ha promovido la \textit{bioeconomía} como paradigma que integra sostenibilidad ecológica, producción basada en recursos biológicos y nuevas formas de organización territorial [7]. 
Red Biet conecta estos conceptos: bioeconomía, identidad digital, blockchain y gobernanza descentralizada.

\section{Motivación de Red Biet}
La motivación central es pasar de una lógica de “beneficiarios pasivos” a una red de Biets como unidades productivas, sociales y educativas que:
\begin{itemize}
    \item participan activamente en la gobernanza del uso de recursos,
    \item generan y registran valor (económico, social, ambiental),
    \item acceden a incentivos y oportunidades de mercado de forma transparente.
\end{itemize}

%----------------------------------
\chapter{Diseño conceptual: Red Biet y Biet Coin (BGT)}
%----------------------------------

\section{Red Biet como red de unidades vivas}
Red Biet se concibe como una malla de Biets (unidades vivas) conectadas por contratos inteligentes. 
Cada Biet puede representar:
\begin{itemize}
    \item una unidad doméstica,
    \item una organización comunitaria,
    \item un grupo productivo,
    \item un laboratorio social o educativo.
\end{itemize}

Cada Biet cuenta con:
\begin{itemize}
    \item Identidad DID.
    \item Perfil productivo y social.
    \item Historial de interacciones on-chain.
    \item Reputación asociada a indicadores objetivos.
\end{itemize}

\section{Biet Coin (BGT)}
El token Biet Coin (BGT) cumple cuatro funciones:
\begin{enumerate}
    \item \textbf{Gobernanza:} da derecho a voto y a proponer cambios en la DAO.
    \item \textbf{Incentivos:} recompensa producción, impacto social y ambiental.
    \item \textbf{Reputación:} refleja trayectoria de contribución del Biet.
    \item \textbf{Medio de intercambio:} facilita intercambios en mercados locales y en el marketplace Web3 de la red.
\end{enumerate}

La emisión de BGT no se basa en minería tradicional, sino en \textit{pruebas de impacto} y contribuciones verificables, alineadas con objetivos de desarrollo.

%----------------------------------
\chapter{Arquitectura Web3 e integración técnica}
%----------------------------------

\section{Capas de la arquitectura}
La arquitectura de Red Biet se organiza en capas:

\begin{itemize}
    \item \textbf{Capa de identidad (DID):} basada en el estándar W3C DID [4], con anclajes on-chain.
    \item \textbf{Capa de contratos inteligentes (EVM):} incluye BGT, BietRegistry, BietDAO, Treasury, BietVault y otros módulos.
    \item \textbf{Capa de aplicación Web:} desarrollada en Next.js, integrando WalletConnect, Metamask u otros proveedores.
    \item \textbf{Capa de datos y analítica:} paneles para monitoreo de impacto social y trazabilidad financiera.
\end{itemize}

\section{Integración con Next.js y EVM}
La dApp de Red Biet se desarrolla con Next.js, conectándose a la red EVM (por ejemplo, Base testnet) mediante librerías como \texttt{ethers.js} o \texttt{wagmi}. 
El flujo típico es:
\begin{enumerate}
    \item El usuario se autentica con su wallet y DID.
    \item La dApp consulta los contratos (BietRegistry, BietDAO, BGT) para mostrar estado.
    \item Las interacciones (votos, propuestas, transacciones) se firman desde el cliente y se envían a la red.
    \item Los eventos on-chain alimentan dashboards y paneles de impacto.
\end{enumerate}

\section*{Arquitectura del sistema}
\addcontentsline{toc}{section}{Arquitectura del sistema}
\par\noindent
\begin{otherlanguage*}{english}
\begin{figure}[H]
\centering
\begin{tikzspanish}[node distance=10mm, >=Latex]
\tikzset{
 layer/.style={rectangle,draw,rounded corners,align=center,minimum width=40mm,minimum height=8mm,fill=white},
 box/.style={rectangle,draw,rounded corners,align=center,minimum width=28mm,minimum height=6mm,fill=white},
 every node/.style={font=\small}
}

\node[layer] (ui) {Capa de Aplicación Web\\(Next.js, Dashboard, UX comunitaria)};
\node[layer,below=of ui] (logic) {Capa de Lógica Web3\\(APIs, WalletConnect, DIDKit)};
\node[layer,below=of logic] (contracts) {Capa de Contratos Inteligentes\\(EVM L2 -- Base u otra compatible)};
\node[layer,below=of contracts] (infra) {Infraestructura Blockchain\\Nodos, indexadores, oráculos};

\node[box,right=18mm of contracts] (bgt) {BGT Token};
\node[box,above=3mm of bgt] (dao) {BietDAO};
\node[box,below=3mm of bgt] (vault) {BietVault};
\node[box,right=18mm of dao] (reg) {BietRegistry};

\draw[->] (ui) -- (logic);
\draw[->] (logic) -- (contracts);
\draw[->] (contracts) -- (infra);

\draw[->] (contracts.east) -- (bgt.west);
\draw[->] (bgt.north) -- (dao.south);
\draw[->] (bgt.south) -- (vault.north);
\draw[->] (dao.east) -- (reg.west);

\node[box,above=8mm of ui] (users) {Biets, organizaciones, veedores};
\draw[->] (users) -- (ui);

\end{tikzspanish}
\caption{Arquitectura general Web3 de la Red Biet.}
\end{figure}
\end{otherlanguage*}

%----------------------------------
\chapter{Flujos operativos y financieros}
%----------------------------------

\section{Flujo operativo de un Biet}
El flujo operativo recoge las fases que atraviesa una comunidad para convertirse en un conjunto de Biets productivos integrados en la red.

\subsection*{Flujo operativo}
\addcontentsline{toc}{subsection}{Flujo operativo}
\par\noindent
\begin{otherlanguage*}{english}
\begin{figure}[H]
\centering
\begin{tikzspanish}[
  node distance=10mm,
  every node/.style={font=\small},
  box/.style={rectangle,draw,rounded corners,align=center,minimum width=35mm,minimum height=8mm,fill=white},
  arrow/.style={-{Latex[length=3mm]},thick}
]
  % First row
  \node[box] (org) {Organización comunitaria};
  \node[box,right=of org] (diag) {Diagnóstico y selección de Biets};
  \node[box,right=of diag] (form) {Formación integral (BietLab)};

  % Second row
  \node[box,below=15mm of diag] (create) {Creación y formalización de Biets};
  \node[box,right=of create] (prod) {Producción y servicios en red};
  \node[box,right=of prod] (eval) {Evaluación e incentivos BGT};

  % Flow
  \draw[arrow] (org) -- (diag);
  \draw[arrow] (diag) -- (form);
  \draw[arrow] (form.south) -- (create);
  \draw[arrow] (create) -- (prod);
  \draw[arrow] (prod) -- (eval);

  % Feedback
  \draw[arrow,green!60!black] (eval.north) .. controls +(0,15mm) and +(0,15mm) .. (form.north);
  \draw[arrow,green!60!black] (prod.west) .. controls +(-22mm,0) and +(22mm,0) .. (diag.south);

\end{tikzspanish}
\caption{Flujo operativo para la creación y consolidación de Biets.}
\end{figure}
\end{otherlanguage*}

\section{Flujo de dinero en la red}

\subsection*{Flujo financiero}
\addcontentsline{toc}{subsection}{Flujo financiero}
\par\noindent
\begin{otherlanguage*}{english}
\begin{figure}[H]
\centering
\begin{tikzspanish}[
    node distance=24mm,
    every node/.style={font=\footnotesize,align=center,inner sep=4pt},
    fund/.style={ellipse, draw, align=center, minimum width=36mm, minimum height=10mm, fill=white},
    entity/.style={rectangle, draw, rounded corners, align=center, minimum width=38mm, minimum height=10mm, fill=white},
    arrow/.style={-{Latex[length=3mm]}, thick},
    flowlabel/.style={font=\scriptsize, fill=white, inner sep=1pt}
]

% NODES
\node[fund] (donor) {Donantes / Organismos\\Internacionales};
\node[entity, right=30mm of donor] (platform) {Fondo Red Biet\\(Tesorería DAO)};
\node[entity, right=30mm of platform] (bgtfund) {Reserva BGT\\(Liquidez + Incentivos)};

\node[entity, below=22mm of platform] (biets) {Biets Productivos\\(Unidades sociales)};
\node[entity, right=30mm of biets] (markets) {Mercado local\\y Marketplace Web3};

\node[fund, above=22mm of platform] (cit) {Veeduría Ciudadana\\Auditoría comunitaria};

% ARROWS
\draw[arrow] (donor) -- node[flowlabel,above]{Subvenciones / fondos sociales} (platform);

\draw[arrow] (platform) -- node[flowlabel,left]{Asignación automática\\(smart contracts)} (biets);
\draw[arrow] (platform) -- node[flowlabel,above]{Provisión / compra BGT} (bgtfund);

\draw[arrow] (biets) -- node[flowlabel,below]{Productos / servicios} (markets);
\draw[arrow] (markets) -- node[flowlabel,below]{Ingresos FIAT / cripto} (biets);

\draw[arrow] (biets) -- node[flowlabel,right]{Reinversión productiva} (platform);
\draw[arrow] (bgtfund) -- node[flowlabel,right]{Recompensas BGT} (biets);

\draw[arrow,dashed] (cit) -- node[flowlabel,right]{Auditoría pública on-chain} (platform);

\end{tikzspanish}
\caption{Flujo financiero corregido y centrado en la Red Biet.}
\end{figure}
\end{otherlanguage*}

%----------------------------------
\chapter{Gobernanza DAO, tokenomics y auditoría}
%----------------------------------

\section{Gobernanza DAO}
La DAO de Red Biet define reglas para:
\begin{itemize}
    \item presentar propuestas,
    \item priorizar proyectos,
    \item asignar presupuestos y BGT,
    \item modificar parámetros del protocolo.
\end{itemize}

\subsection*{Flujo de gobernanza}
\addcontentsline{toc}{subsection}{Flujo de gobernanza}
\par\noindent
\begin{otherlanguage*}{english}
\begin{figure}[H]
\centering
\begin{tikzspanish}[
  node distance=12mm,
  every node/.style={font=\small},
  box/.style={rectangle,draw,rounded corners,align=center,minimum width=32mm,minimum height=8mm,fill=white},
  arrow/.style={-{Latex[length=3mm]},thick}
]
  \node[box] (prop) {Propuesta comunitaria};
  \node[box,right=of prop] (review) {Revisión técnica / social};
  \node[box,right=of review] (vote) {Votación DAO\\(ponderada por BGT y reputación)};
  \node[box,right=of vote] (exec) {Ejecución on-chain\\(smart contract)};
  \node[box,below=of vote] (report) {Reporte y monitoreo};
  \node[box,below=of prop] (audit) {Auditoría y evaluación};

  \draw[arrow] (prop) -- (review);
  \draw[arrow] (review) -- (vote);
  \draw[arrow] (vote) -- (exec);
  \draw[arrow] (exec) -- (report);
  \draw[arrow] (report) -- (audit);
  \draw[arrow] (audit.west) .. controls +(-22mm,0) and +(0,-15mm) .. (prop.south);

\end{tikzspanish}
\caption{Flujo de gobernanza DAO en Red Biet.}
\end{figure}
\end{otherlanguage*}

\section{Tokenomics y modelo criptográfico}
La tokenomics de BGT se diseña para alinear incentivos con objetivos de desarrollo:

\begin{itemize}
    \item \textbf{Suministro inicial:} fijado por la DAO en el despliegue.
    \item \textbf{Emisión dinámica:} ligada a indicadores de impacto social y productivo.
    \item \textbf{Pools de incentivos y liquidez:} gestionados por el Treasury.
    \item \textbf{Modelo criptográfico:} uso de firmas ECDSA/ECDSA-secp256k1 (como en Ethereum) y gestión de roles mediante contratos de acceso \textit{role-based}.
\end{itemize}

\subsection*{Tokenomics BGT}
\addcontentsline{toc}{subsection}{Tokenomics BGT}
\par\noindent
\begin{otherlanguage*}{english}
\begin{figure}[H]
\centering
\begin{tikzspanish}[
  node distance=12mm,
  every node/.style={font=\small},
  box/.style={rectangle,draw,rounded corners,align=center,minimum width=34mm,minimum height=8mm,fill=white},
  arrow/.style={-{Latex[length=3mm]},thick}
]
  \node[box] (treasury) {Treasury DAO};
  \node[box,above left=of treasury] (poolIncent) {Pool de incentivos BGT};
  \node[box,above right=of treasury] (poolLiq) {Pool de liquidez / DEX};
  \node[box,below left=of treasury] (biets) {Biets productivos};
  \node[box,below right=of treasury] (partners) {Aliados / Programas sociales};

  \draw[arrow] (treasury) -- node[left]{Asignación} (poolIncent);
  \draw[arrow] (treasury) -- node[right]{Provisión} (poolLiq);
  \draw[arrow] (poolIncent) -- node[left]{Recompensas} (biets);
  \draw[arrow] (biets) -- node[below]{Datos de impacto / méritos} (treasury);
  \draw[arrow] (poolLiq) -- node[right]{Swap / mercado} (partners);
  \draw[arrow] (partners) -- node[below]{Aportes financieros} (treasury);

\end{tikzspanish}
\caption{Esquema simplificado de la tokenomics de BGT.}
\end{figure}
\end{otherlanguage*}

\section{Modelo de auditoría y transparencia}
La auditoría en Red Biet combina:
\begin{itemize}
    \item registros on-chain accesibles vía exploradores y APIs,
    \item dashboards públicos en la dApp,
    \item participación de veedurías ciudadanas,
    \item acompañamiento de universidades y ONGs.
\end{itemize}

\subsection*{Modelo de auditoría}
\addcontentsline{toc}{subsection}{Modelo de auditoría}
\par\noindent
\begin{otherlanguage*}{english}
\begin{figure}[H]
\centering
\begin{tikzspanish}[
  node distance=10mm,
  every node/.style={font=\small},
  box/.style={rectangle,draw,rounded corners,align=center,minimum width=36mm,minimum height=8mm,fill=white},
  arrow/.style={-{Latex[length=3mm]},thick}
]
  \node[box] (data) {Datos on-chain (transacciones, votos, emisión BGT)};
  \node[box,below=of data] (dash) {Dashboards públicos y analítica};
  \node[box,below left=of dash] (cit) {Veeduría ciudadana};
  \node[box,below right=of dash] (acad) {Universidades / ONGs};
  \node[box,below=of dash,yshift=-20mm] (dao) {DAO Red Biet};

  \draw[arrow] (data) -- (dash);
  \draw[arrow] (dash) -- (cit);
  \draw[arrow] (dash) -- (acad);
  \draw[arrow] (cit) -- node[below]{Alertas / propuestas} (dao.west);
  \draw[arrow] (acad) -- node[below]{Informes / métricas} (dao.east);
  \draw[arrow] (dao.north) -- node[right]{Decisiones / cambios de protocolo} (data.east);

\end{tikzspanish}
\caption{Modelo de auditoría participativa y transparencia en Red Biet.}
\end{figure}
\end{otherlanguage*}

%----------------------------------
\chapter{Impacto social, ODS e indicadores}
%----------------------------------

\section{Alineación con los ODS}
Red Biet se alinea explícitamente con los siguientes Objetivos de Desarrollo Sostenible (ODS):
\begin{itemize}
    \item \textbf{ODS 1:} Fin de la pobreza.
    \item \textbf{ODS 8:} Trabajo decente y crecimiento económico.
    \item \textbf{ODS 10:} Reducción de las desigualdades.
    \item \textbf{ODS 16:} Paz, justicia e instituciones sólidas.
\end{itemize}

\section{Indicadores propuestos}
Algunos indicadores para evaluar impacto:
\begin{itemize}
    \item Número de Biets activos y sostenibles en el tiempo.
    \item Incremento del ingreso promedio por Biet.
    \item Porcentaje de transacciones auditables públicamente.
    \item Participación efectiva en votaciones DAO.
    \item Reducción de dependencia de subsidios puramente asistenciales.
\end{itemize}

%----------------------------------
\chapter{Anexo técnico: esbozo de contratos}
%----------------------------------

\section{Ejemplo (esquemático) — BGT token (Solidity)}
\begin{verbatim}
pragma solidity ^0.8.19;
import "@openzeppelin/contracts/token/ERC20/ERC20.sol";

contract BGT is ERC20 {
  constructor(uint256 initialSupply) ERC20("Biet Coin", "BGT") {
    _mint(msg.sender, initialSupply);
  }
}
\end{verbatim}

\section{Otros contratos (esbozo conceptual)}
De manera análoga se implementan:
\begin{itemize}
    \item \textbf{BietRegistry:} registro de Biets, roles y estados.
    \item \textbf{BietDAO:} módulo de propuestas, votos y ejecución.
    \item \textbf{Treasury:} gestión de fondos, reservas y distribución.
    \item \textbf{BietVault:} almacenamiento de reputación e indicadores.
\end{itemize}

Las firmas se basan en el esquema ECDSA de Ethereum; el control de acceso puede implementarse con \texttt{AccessControl} de OpenZeppelin.

%----------------------------------
\chapter*{Referencias}
\addcontentsline{toc}{chapter}{Referencias}
%----------------------------------


\begin{thebibliography}{99}

\bibitem{OECD2019}
OECD (2019).
\newblock Social Expenditure Database (SOCX).
\newblock \href{https://www.oecd.org/social/expenditure.htm}{https://www.oecd.org/social/expenditure.htm}.

\bibitem{WorldBank2020}
World Bank (2020).
\newblock World Development Report 2020.
\newblock \href{https://www.worldbank.org}{https://www.worldbank.org}.

\bibitem{TI2021}
Transparency International (2021).
\newblock Corruption Perceptions Index 2020.
\newblock \href{https://www.transparency.org/en/cpi/2020}{https://www.transparency.org/en/cpi/2020}.

\bibitem{Allen2017}
Allen, C. et al. (2017).
\newblock Decentralized Identifiers (DIDs) v1.0.
\newblock W3C Recommendation.
\newblock \href{https://www.w3.org/TR/did-core/}{https://www.w3.org/TR/did-core/}.

\bibitem{Buterin2014}
Buterin, V. (2014).
\newblock DAOs, DACs, DAs and More: An Incomplete Terminology Guide.
\newblock Ethereum Blog.
\newblock \href{https://blog.ethereum.org}{https://blog.ethereum.org}.

\bibitem{Cong2019}
Cong, L. W., \& He, Z. (2019).
\newblock Blockchain Disruption and Smart Contracts.
\newblock \textit{Review of Financial Studies}.
\newblock \href{https://doi.org/10.1093/rfs/hhz007}{https://doi.org/10.1093/rfs/hhz007}.

\bibitem{EC2018}
European Commission (2018).
\newblock A sustainable bioeconomy for Europe.
\newblock \href{https://ec.europa.eu}{https://ec.europa.eu}.

\bibitem{Putnam2000}
Putnam, R. D. (2000).
\newblock Bowling Alone: The Collapse and Revival of American Community.
\newblock Simon \& Schuster.

\bibitem{Ravallion1995}
Ravallion, M. (1995).
\newblock Measuring Poverty.
\newblock \textit{The World Bank Economic Review}, 9(3).

\bibitem{Sen1999}
Sen, A. (1999).
\newblock Development as Freedom.
\newblock Oxford University Press.

\end{thebibliography}

\end{document}